\documentclass[10pt,a4paper]{amsart}
\usepackage[margin=19mm]{geometry}
\usepackage{amsmath,amssymb,mathtools}
\usepackage[scr=rsfs,cal=euler]{mathalfa}
\usepackage{xfrac}
\usepackage[unicode,hidelinks,bookmarks=false]{hyperref}
\newcommand{\ie}{i.e.\ }
\newcommand{\cf}{cf.\ }
\newcommand{\resp}{respectively\ }
\newcommand{\Subsection}{\subsection}
\providecommand{\nobreakdash}{\nobreak\hbox{-}}
\setlength{\emergencystretch}{2em}
\multlinegap0pt
\usepackage{bm}
\usepackage{bbm}
\usepackage{dsfont}
\usepackage{xspace}
\usepackage{cancel}
\usepackage{mathtools}
\usepackage{amsthm}
\makeatletter
\@ifundefined{thm}{\newtheorem{thm}{Thm}}{}
\@ifundefined{ex}{\newtheorem{ex}[thm]{Example}}{}
\makeatother
\title[Maximal dilatation on nonorientable surfaces]{Maximal dilatation on nonorientable surfaces}
\author{Ji-Young Ham, Joongul Lee}
\date{Source: arXiv:2508.08323v1 (CC BY 4.0)}
\begin{document}
\maketitle

\noindent\emph{Source and licence.} Maximal dilatation on nonorientable surfaces. Version arXiv:2508.08323v1 (CC BY 4.0), distributed under the Creative Commons Attribution 4.0 International licence, as recorded in the arXiv licence metadata for this exact version and shown on the abstract page https://arxiv.org/abs/2508.08323v1. Full rights notice: licenses/arxiv-2508.08323-CC-BY-4.0.txt.\par\smallskip

\noindent\textit{Editorial scope.} Counted unit: the Ex whose source label is \texttt{ex:2} in the article (source0.tex lines 330--332, source label \texttt{ex:2}), delivered together with the article's own complete original proof of it (source0.tex lines 334--413, one \texttt{proof} environment of 80 lines). No mathematics is written, repaired or re-derived: statement and proof are the article's own text line for line. Required context retained uncounted: none. Every definition, display and label of those lines is retained unchanged. Omitted from this package and not redistributed: 1--329, 333--333, 414--675. Nothing inside a retained excerpt is omitted or altered: every retained body is delivered exactly as the article's lines are written, so no omission receipt of any kind accompanies this package. Citations: the retained excerpts carry no citation command at all, so no bibliography entry of the article is transcribed and no reference is redistributed. Reference adaptation: none was needed and none was made; every reference inside the delivered unit resolves inside it, so no unresolved reference or citation is produced. Printed slips kept literal: no printed word, symbol, hypothesis or number of the counted statement or of its proof was altered. Layout and class: the article's own class is \texttt{amsart} with packages including amsmath, amsfonts, amssymb, mathabx, graphicx, url, hyperref, comment, makecell, boldline, hhline, colortbl, xcolor, textcomp; the wrapper is the lane's pinned \texttt{amsart} editorial wrapper at 10pt on A4, which restates the theorem-like environments the retained text needs and adds the article's own macro definitions that the retained text needs (none), with extra wrapper packages (bm, bbm, dsfont, xspace, cancel, mathtools). The delivered numbering is the unit's own. Assets: no retained span contains a figure environment, a graphics-inclusion command, a PDF-page inclusion, a table or a TikZ picture; no image file is copied into this package, the package declares no asset dependency and no image file is redistributed with it. Licence: version arXiv:2508.08323v1 of this article is distributed under the Creative Commons Attribution 4.0 International licence, as recorded in the arXiv licence metadata for this exact version and shown on the abstract page https://arxiv.org/abs/2508.08323v1. A full rights notice is delivered at licenses/arxiv-2508.08323-CC-BY-4.0.txt. Characters: every retained excerpt is pure ASCII, so the delivered engine, XeLaTeX with its default Latin Modern text font, reports no missing character.
\begin{ex} \label{ex:2}
$x^5-x^3-x^2-1$ is irreducible.
\end{ex} 

\begin{proof}
$\bullet$ $x^5-x^3-x^2-1$ does not have linear factors.
\begin{enumerate} 
\item[I.] Suppose $x^5-x^3-x^2-1=(x^4+a_3 x^3+a_2 x^2+a_1 x+1)(b_1 x+b_0)$ where $b_0^2=1$ and $b_1^2=1$. 
Then, by comparing the coefficients, we have

%%%%%%%%%%%%%%%%%%%%
\begin{enumerate}
\item[(0)] $b_0=-1$ (the constants), \\
\item[(1)] $b_1+a_1 b_0=0$ (the coefficients of $x^1$), \\
\item[(2)] $a_1 b_1+a_2 b_0=-1$ (the coefficients of $x^2$) $\quad \Rightarrow \quad a_1 b_1+a_2 b_0+1=0$,\\
\item[(3)] $a_2 b_1+a_3 b_0=-1$ (the coefficients of $x^3$) $\quad \Rightarrow \quad a_2 b_1+a_3 b_0+1=0$, \\
\item[(4)] $a_3 b_1+b_0=0$ (the coefficients of $x^4$), \\
\item[(5)] $b_1=1$.
\end{enumerate}


We compute the resultant symmetrically in such a way that when we take the resultant of $(i)$ and $(j)$ with respect to $a_{t}$, we also take the resultant of $(5-i)$ and $(5-j)$ with respect to $a_{4-t}$.

\begin{enumerate}
\item[(6)] $res_{a_1}((1),(2))=-b_1^2+a_2b_0^2+b_0=-1+a_2+b_0$,
\item[(7)] $res_{a_3}((4),(3))=a_2 b_1^2+b_1-b_0^2=a_2+b_1-1$,
\end{enumerate}

Now,

\begin{equation*}
res_{a_2}((6),(7))=b_1-b_0
\end{equation*}

But then, since $b_1=1$, $b_0=1$ which is inconsistent with $(0)$.
%%%%%%%%%%%%%%%%%%%%%%%%%%%%%%%%%%%%%%%%%%%%%%%%%%
\item[II.] Suppose $x^5-x^3-x^2-1=(a_4x^4+a_3 x^3+a_2 x^2+a_1 x+a_0)(x+1)$ where $a_0^2=1$ and $a_4^2=1$. 
Then, by comparing the coefficients, we have

%%%%%%%%%%%%%%%%%%%%
\begin{enumerate}
\item[(0)] $a_0=-1$ (the constants), \\
\item[(1)] $a_1+ a_0=0$ (the coefficients of $x^1$), \\
\item[(2)] $a_1+a_2=-1$ (the coefficients of $x^2$) $\quad \Rightarrow \quad a_1+a_2+1=0$,\\
\item[(3)] $a_2+a_3=-1$ (the coefficients of $x^3$) $\quad \Rightarrow \quad a_2+a_3+1=0$, \\
\item[(4)] $a_3+a_4=0$ (the coefficients of $x^4$), \\
\item[(5)] $a_4=1$.
\end{enumerate}


We compute the resultant symmetrically in such a way that when we take the resultant of $(i)$ and $(j)$ with respect to $a_{t}$, we also take the resultant of $(5-i)$ and $(5-j)$ with respect to $a_{4-t}$.

\begin{enumerate}
\item[(6)] $res_{a_1}((1),(2))=a_2-a_0+1$,
\item[(7)] $res_{a_3}((4),(3))=a_2-a_4+1$.
\end{enumerate}

Now,

\begin{equation*}
res_{a_2}((6),(7))=a_0-a_4
\end{equation*}

But then, since $a_4=1$, $a_0=1$ which is inconsistent with $(0)$.

\end{enumerate} 
%%%%%%%%%%%%%%%%%%%%%%%%%%%%%%%%%%%%%%%%%%%%%%%%%%%%%%

$\bullet$ $x^5-x^3-x^2-1$ does not have quadratic factors.
\begin{enumerate} 
\item[I.] Suppose $x^5-x^3-x^2-1=(x^3+a_2 x^2+a_1 x+1)(b_2 x^2+b_1 x+b_0)$ where $b_0^2=1$ and $b_2^2=1$.
Similarly, we symmetrically take the resultant several times until we get two equations with three variables $b_2$,$b_1$, and $b_0$.
Now, taking the resultant one more time with respect to $b_1$, we get
 \begin{equation*}
 4(1-b_0 b_2)=4b_2(b_2-b_0) .
\end{equation*}
But then, the system becomes inconsistent.
\item[II.] Suppose $x^5-x^3-x^2-1=(a_3 x^3+a_2 x^2+a_1 x+a_0)(x^2+b_1 x+1)$.
Similarly, we have $2(1-a_0 a_3)=2a_3(a_3-a_0).$
But then, the system becomes inconsistent.
\end{enumerate}

Therefore, by two bullets, $x^5-x^3-x^2-1$ is irreducible.
\end{proof}

\end{document}
